% TOP_PROBLEM: {"id": 30006677, "title": "Extended Riemann hypothesis for Dedekind zeta functions", "category_id": 1, "category": {"id": 1, "name": "number_theory", "display_name": "Number Theory", "description": "Properties of integers, prime numbers, Diophantine equations.", "slug": "number-theory", "order_index": 1, "created_at": "2026-07-31T15:26:25.670Z"}, "status": "open"}
% FIRST_POSTED: null
% !TeX program = lualatex
% Compile twice: lualatex open_problems_500.tex
% All references and prose are embedded; no BibTeX database or external inputs.
\documentclass[11pt,a4paper]{article}
\usepackage[margin=24mm,headheight=14pt]{geometry}
\usepackage{fontspec}
\setmainfont{Latin Modern Roman}
\setsansfont{Latin Modern Sans}
\setmonofont{Latin Modern Mono}
\usepackage{amsmath,amssymb,mathtools}
\usepackage{unicode-math}
\setmathfont{Latin Modern Math}
\usepackage{microtype}
\usepackage{enumitem,xurl,needspace}
\usepackage[dvipsnames]{xcolor}
\usepackage{hyperref}
\hypersetup{unicode=true,colorlinks=true,linkcolor=MidnightBlue,urlcolor=MidnightBlue,
 pdftitle={Mathematical Research Notebook},pdfauthor={Source-annotated compilation},bookmarksnumbered=true}
\usepackage{bookmark}
\usepackage{tocloft}
\setlength{\cftsecnumwidth}{3em}
\cftsetpnumwidth{2.5em}
\cftsetrmarg{4em}
\usepackage{titlesec}
\titleformat{\section}{\Large\bfseries\raggedright}{\thesection}{0.7em}{}
\usepackage{fancyhdr}
\pagestyle{fancy}\fancyhf{}
\fancyhead[L]{\small\sffamily Open Problems Math | Research notebook}
\fancyhead[R]{\small 27 September 2026}
\fancyfoot[C]{\thepage}
\setlength{\parindent}{0pt}
\setlength{\parskip}{4pt plus 1pt}
\setlength{\emergencystretch}{3em}
\setcounter{tocdepth}{1}
\setcounter{secnumdepth}{1}
\setlist[itemize]{leftmargin=3em,itemsep=3pt,topsep=4pt,parsep=0pt}
\allowdisplaybreaks
\newcommand{\N}{\mathbb{N}}
\newcommand{\Z}{\mathbb{Z}}
\newcommand{\Q}{\mathbb{Q}}
\newcommand{\R}{\mathbb{R}}
\newcommand{\C}{\mathbb{C}}
\newcommand{\F}{\mathbb{F}}
\newcommand{\A}{\mathbb{A}}
\newcommand{\PP}{\mathbb P}
\newcommand{\E}{\mathbb{E}}
\newcommand{\eps}{\varepsilon}
\newcommand{\dd}{\,\mathrm d}
\newcommand{\sref}[2]{\hyperlink{source:#1:#2}{[S#2]}}
\newcommand{\eref}[2]{\hyperlink{extra:#1:#2}{[E#2]}}
% Turn the next switch off for a shorter reading copy. The quotations preserve
% every exactTarget field, including qualifications omitted from a short summary.
\newif\ifshowcatalogue
\showcataloguetrue
\newcommand{\cataloguescope}[1]{\ifshowcatalogue
 \paragraph{Exact scope in the supplied catalogue.}
 \begin{quote}\small #1\end{quote}\fi}
\usepackage{etoolbox}
\AtBeginEnvironment{itemize}{\small}
% Standalone entry; original section content is delimited for verification.
\begin{document}
\setcounter{section}{0}
%% BEGIN ORIGINAL SECTION CONTAINER
% ================================================================
% problemId: 30006677
% Canonical problem ID: 30006677
% exactTarget SHA-256: 48d6696cd553434cea27f3244febc5017e4fd62fbac401b6cc39aa55623e4ebf
\clearpage
\section*{\texorpdfstring{Extended Riemann hypothesis for Dedekind zeta functions}{Extended Riemann hypothesis for Dedekind zeta functions}}\label{30006677}
\subsection{Definitions and mathematical statement}
For a number field $K$, let $\mathcal O_K$ be its ring of integers and $N\mathfrak a=\#(\mathcal O_K/\mathfrak a)$ the norm of a nonzero integral ideal. Define
\[
 \zeta_K(s)=\sum_{0\ne\mathfrak a\subseteq\mathcal O_K}(N\mathfrak a)^{-s}
 \qquad(\Re s>1),
\]
and continue meromorphically. The assertion is $\zeta_K(\rho)=0$ and $0<\Re\rho<1\Rightarrow\Re\rho=1/2$, for every finite extension $K/\Q$. The rational field itself is included.
\par\smallskip\textit{Formulation and concept references:} \sref{30006677}{1}.
\cataloguescope{For every algebraic number field K (a finite extension of Q, including K=Q), does every zero s of its Dedekind zeta function \ensuremath{\zeta}\_K(s) in the critical strip 0<Re(s)<1 satisfy Re(s)=1/2?}
\subsection{Short English statement}
For every number field, must the zeros of its ideal-counting zeta function inside the critical strip lie halfway across it?
%% BEGIN RESEARCH INSERTION
\subsection{Research attempt}
\paragraph{Current frontier.}
\textit{Research check: 27 September 2026.}
The target includes every number field, including $\Q$.
\sref{30006677}{1}. Explicit prime-ideal estimates under the corresponding
Riemann hypothesis are available, for example Grenie--Molteni,
Theorem 1.1; the hypothesis in that theorem must not be removed when
using its conclusion. Their Section 2 records the completed zeta
function and functional equation. \eref{30006677}{1}.
The titles of the inherited purported proofs are not adopted as
verification. This attempt instead isolates an unconditional analytic
estimate whose proof would settle the target, and tests whether
positivity of ideal-counting coefficients can force it.

\paragraph{Attack route.}
Three possible routes are a positivity criterion for the completed
function, a prime-ideal error estimate, or passage from abelian to
nonabelian fields. A positivity argument needs information stronger
than nonnegative Euler coefficients. The error-estimate route needs
cancellation down to the square-root scale. The field-extension route
needs analytic control beyond products of Dirichlet $L$-functions.
We pursue the second, while calculating exactly what the third gives
for cyclotomic fields.

\paragraph{Attempt: coefficient bounds with all ramification retained.}
For a prime ideal $\mathfrak p$, define
$\Lambda_K(\mathfrak p^r)=\log N\mathfrak p$ for $r\geq1$, and
zero on other ideals. Group terms according to their integer norms:
\[
 -\frac{\zeta_K'}{\zeta_K}(s)
 =\sum_{n\geq1}\frac{b_K(n)}{n^s},\qquad
 \psi_K(x)=\sum_{n\leq x}b_K(n),\qquad \Re s>1. \tag{10.1}
\]
Writing $d=[K:\Q]$, if the residue degrees of primes over $p$ are
$f_1,\ldots,f_g$, then
\[
 b_K(p^m)=\Bigl(\sum_{j:f_j\mid m}f_j\Bigr)\log p,
 \qquad 0\leq b_K(p^m)\leq d\log p. \tag{10.2}
\]
At a ramified prime the same formula holds: ramification indices enter
$\sum_j e_jf_j=d$ and give $\sum_jf_j\leq d$, not an omitted correction
to (10.2). Non-prime-power coefficients vanish. Absolute convergence of
the logarithmic derivative in $\Re s>1$ follows by comparison with
$d\sum_n\Lambda(n)n^{-\sigma}$. This justifies differentiating the
Euler product there. The definitions are the prime-ideal normalization
in \eref{30006677}{1}, Section 1; (10.2) is the explicit local calculation.

\paragraph{Attempt: a sufficient estimate and a complete Mellin argument.}
Suppose that for every fixed field $K$ and every $\varepsilon>0$,
\[
 \psi_K(x)=x+O_{K,\varepsilon}(x^{1/2+\varepsilon})
 \quad(x\geq2). \tag{10.3}
\]
Then the exact assertion in this entry follows. No constant uniform
in degree or discriminant is needed by this reduction.
For $\Re s>1$, summation by parts, using (10.2) for absolute convergence,
gives
\[
 -\frac{\zeta_K'}{\zeta_K}(s)
 =s\int_1^\infty\psi_K(x)x^{-s-1}\,dx
 =\frac{s}{s-1}+s\int_1^\infty E_K(x)x^{-s-1}\,dx,
 \qquad E_K(x)=\psi_K(x)-x. \tag{10.4}
\]
There is no discarded lower-end constant: $\psi_K(1)=0$ and
$E_K(x)$ is defined on the whole interval $[1,\infty)$.

Let $C$ be a compact subset of $\Re s>1/2$, and let
$\sigma_0=\min_{s\in C}\Re s>1/2$. Choose
$0<\varepsilon<\sigma_0-1/2$. The last integrand is bounded, uniformly
on $C$, by a constant times
$x^{-1-(\sigma_0-1/2-\varepsilon)}$ for $x\geq2$.
Its derivatives in $s$ introduce only powers of $\log x$, still
integrable. The integral therefore defines a holomorphic function on
$\Re s>1/2$. By uniqueness of meromorphic continuation, the right side
of (10.4) equals the logarithmic derivative throughout that half-plane.
It has no pole there except at $s=1$.

If $\rho\ne1$ in that half-plane were a zero of $\zeta_K$ of multiplicity
$m\geq1$, its negative logarithmic derivative would have a pole of
residue $-m$, a contradiction. The completed functional equation
$\xi_K(s)=\xi_K(1-s)$, with gamma factors nonzero and finite in the
open strip, reflects any left-half-strip zero to a right-half-strip
zero. \eref{30006677}{1}, equations (2.1)--(2.3).
Thus all open-strip zeros have real part $1/2$. The known direction
from the hypothesis to estimates of the shape (10.3) is not used to
establish (10.3); doing that would be circular.

\paragraph{Attempt: exactly how much the abelian reduction buys.}
Take $K=\Q(\zeta_q)$, first with $q>1$, and let
$G=(\Z/q\Z)^\times$. For $p\nmid q$, its residue degree is
$f=\operatorname{ord}_G(p)$ and it has $|G|/f$ primes above it.
The character values $\chi(p)$, as $\chi$ ranges over $\widehat G$,
contain each $f$th root of unity with multiplicity $|G|/f$. Hence
\[
 \prod_{\chi\bmod q}(1-\chi(p)t)
       =(1-t^f)^{|G|/f}. \tag{10.5}
\]
This follows by restricting characters to the cyclic subgroup generated
by $p$: every character of that subgroup has $|G|/f$ extensions to
$G$, as seen by decomposing the finite abelian group into cyclic factors.
Using the splitting description of cyclotomic primes, (10.5) gives,
initially in $\Re s>1$,
\[
 \zeta_K(s)=R_q(s)\prod_{\chi\bmod q}L(s,\chi),\qquad
 R_q(s)=\prod_{p\mid q}\prod_{\mathfrak p\mid p}
              (1-(N\mathfrak p)^{-s})^{-1}. \tag{10.6}
\]
Here the $L$-functions are the \emph{imprimitive modulus-$q$} functions,
so they omit all primes dividing $q$. The factor $R_q$ restores exactly
those Euler factors. It is finite, nonzero, and holomorphic on $\Re s>0$.
The cyclotomic splitting facts used in this calculation are in
\eref{30006677}{2}, the discussion of factorization in cyclotomic fields.

Inside the open strip no factor in the product has a pole. Thus a zero
of any Dirichlet factor is a zero of $\zeta_K$, and if every Dirichlet
factor has only central zeros, so does this cyclotomic zeta function.
For $q=1$ the corresponding statement is simply $\zeta_{\Q}=\zeta$.
This proves the dependency that the present all-field assertion implies
UnsolvedMath problem 30006676, and proves the converse only for this cyclotomic subfamily.
It does not express the zeta function of every nonabelian field as
such a product.

\paragraph{Stress test: positivity leaves room for an off-line pole.}
From (10.2) we get only $\psi_K(x)\leq d\psi_{\Q}(x)$; comparison of
nonnegative coefficients does not control the signed error
$\psi_K(x)-x$. Even positivity plus a leading asymptotic is insufficient
in the following precise analytic model. For $1/2<\beta<1$, put
$B_\beta(x)=x-x^\beta$ for $x\geq1$. Then
$B_\beta(1)=0$, $B_\beta'(x)=1-\beta x^{\beta-1}>0$,
and $B_\beta(x)\sim x$. Nevertheless
\[
 s\int_1^\infty B_\beta(x)x^{-s-1}\,dx
       =\frac{s}{s-1}-\frac{s}{s-\beta}\quad(\Re s>1)
\]
has an additional pole at $\beta$. This is not an Euler product or a
number-field counterexample; it refutes the inference from these two
analytic properties alone. Similarly, a finite verification of zeros
or finitely many prime-ideal counts cannot establish (10.3) at arbitrarily
large heights or for all fields.

\paragraph{Outcome.}
\textbf{Outcome: Conditional reduction.}
The Mellin continuation argument proves that the quantified estimate
(10.3) would settle the exact all-field assertion, with no hidden
uniformity in discriminant. The local computation retains ramification,
and the cyclotomic calculation makes the dependency on Dirichlet GRH
explicit without presuming general Artin holomorphy. No proof of the
required square-root-scale cancellation is obtained. The positivity
model concerns a failed intermediate inference, not the original
conjecture; historical novelty is not asserted.
%% END RESEARCH INSERTION
\subsection{Sources}
\begin{itemize}\raggedright
\item[\textnormal{[S1]}]\hypertarget{source:30006677:1}{} American Institute of Mathematics RH project, The Extended Riemann Hypothesis; undated article18a.
\url{https://www.aimath.org/WWN/rh/articles/html/18a/}.
{\small\textit{Catalogue formulation source.}}
\item[\textnormal{[S2]}]\hypertarget{source:30006677:2}{} A random polynomial with multiplicative coefficients is almost surely irreducible — Varjú and Xu.
\url{https://academic.oup.com/imrn/article/2026/16/rnag178/8761102}.
{\small\textit{Catalogue research/status reference; not a proof endorsement.}}
\item[\textnormal{[S3]}]\hypertarget{source:30006677:3}{} On the Extended, Generalized, and Grand Riemann Hypotheses: A Unified Approach Based on the General Properties of L-Functions — Weicun Zhang, v19.
\url{https://www.preprints.org/manuscript/202506.0481}.
{\small\textit{Catalogue research/status reference; not a proof endorsement.}}
\item[\textnormal{[S4]}]\hypertarget{source:30006677:4}{} Proof of generalized Riemann hypothesis for Dedekind zetas and Dirichlet L-functions — Andrzej Mądrecki.
\url{https://arxiv.org/abs/0706.0256}.
{\small\textit{Catalogue research/status reference; not a proof endorsement.}}
%% BEGIN RESEARCH REFERENCES
\item[\textnormal{[E1]}]\hypertarget{extra:30006677:1}{} Loïc Grenié and Giuseppe Molteni, \emph{Explicit versions of the prime ideal theorem for Dedekind zeta functions under GRH}, Mathematics of Computation 85(298) (2016), 889--906; arXiv:1312.4463v3, Theorem 1.1 and Section 2.
\url{https://arxiv.org/html/1312.4463v3}.
{\small\textit{Added research source. The completed function, functional equation and a conditional prime-ideal estimate; GRH is not dropped. Consulted 27 September 2026.}}
\item[\textnormal{[E2]}]\hypertarget{extra:30006677:2}{} James S. Milne, \emph{Algebraic Number Theory}, course notes version 3.08 (19 July 2020), Chapter 6 and Chapter 8, Example 8.18.
\url{https://www.jmilne.org/math/CourseNotes/ANT.pdf}.
{\small\textit{Added research source. Cyclotomic Frobenius and the residue-degree description for primes not dividing the modulus. Consulted 27 September 2026.}}
%% END RESEARCH REFERENCES
\end{itemize}

%% END ORIGINAL SECTION CONTAINER
\end{document}
